\documentclass[11pt]{article}

\usepackage{amsmath,amsthm,amssymb}
\usepackage{mathtools}
\usepackage[margin=1in]{geometry}
\usepackage[hidelinks]{hyperref}

\newtheorem{theorem}{Theorem}
\newtheorem{lemma}[theorem]{Lemma}
\newtheorem{proposition}[theorem]{Proposition}
\newtheorem{corollary}[theorem]{Corollary}
\newtheorem{definition}[theorem]{Definition}

\theoremstyle{remark}
\newtheorem{remark}[theorem]{Remark}
\newtheorem{example}[theorem]{Example}

\newcommand{\CC}{\mathbb{C}}
\newcommand{\ZZ}{\mathbb{Z}}
\newcommand{\Rep}{\mathrm{Rep}}
\newcommand{\Crys}{\mathrm{Crys}}
\newcommand{\GL}{\mathrm{GL}}
\newcommand{\End}{\mathrm{End}}
\newcommand{\id}{\mathrm{id}}
\newcommand{\Mod}{\text{-mod}}
\DeclareMathOperator{\Hom}{Hom}

\title{The Hecke Transition Algebra Theorem}
\author{Clio}
\date{April 11, 2026}

\begin{document}
\maketitle

\begin{abstract}
We prove that the Hecke algebra $H_q(S_k)$ serves as the \emph{transition algebra}
governing the passage from braided to coboundary to symmetric monoidal structure
on tensor powers of the fundamental representation of $U_q(\mathfrak{sl}_n)$.
The theorem has five parts: (1) Jimbo's classical Schur--Weyl duality (known),
(2) cactus group action at $q=0$ via swap operators factoring through $H_0(S_k)$ (proved here),
(3) the multiplicity bundle decomposition of the coboundary commutor (proved in companion paper),
(4) the Kamnitzer--Tingley coboundary structure lives inside $H_q$ for all $q$ (proved here by combining existing results), and
(5) the algebra chain $H_q \to H_0 \to \CC[S_k]$ mirrors the categorical chain
$\Rep(U_q) \to \Rep(U_0) \to \Crys$ with a character-level identification via Demazure atoms (proved here).
Computational verification is provided for all parts with $n, k \le 5$.
\end{abstract}

\tableofcontents

%% =================================================================
\section{Setup and notation}
%% =================================================================

Let $\mathfrak{g} = \mathfrak{sl}_n$, $V = \CC^n$ the fundamental representation
of the quantum group $U_q(\mathfrak{sl}_n)$ (with $q$ a formal parameter or generic
complex number, $q \ne 0, \pm 1$). Let $H_q(S_k)$ be the Iwahori--Hecke algebra
of type $A_{k-1}$, generated by $T_1, \ldots, T_{k-1}$ subject to:
\begin{align}
  (T_i - q)(T_i + q^{-1}) &= 0, \label{eq:hecke-quad}\\
  T_i T_{i+1} T_i &= T_{i+1} T_i T_{i+1}, \label{eq:braid}\\
  T_i T_j &= T_j T_i \quad (|i-j| \ge 2). \label{eq:commute}
\end{align}

The \emph{cactus group} $J_k$ has generators $\sigma_{[i,j]}$ for $1 \le i < j \le k$
with relations:
\begin{enumerate}
  \item[(C1)] $\sigma_{[i,j]}^2 = e$ \quad (involutivity),
  \item[(C2)] $\sigma_{[i,j]}\, \sigma_{[p,r]}\, \sigma_{[i,j]} = \sigma_{[i+j-r,\, i+j-p]}$
    for $[p,r] \subseteq [i,j]$ \quad (nesting),
  \item[(C3)] $\sigma_{[i,j]}\, \sigma_{[p,r]} = \sigma_{[p,r]}\, \sigma_{[i,j]}$
    for $[i,j] \cap [p,r] = \emptyset$ \quad (disjoint commutativity).
\end{enumerate}

We write $w_0^{[i,j]}$ for the longest element of the parabolic subgroup $S_{\{i,\ldots,j\}}$
(the permutation reversing the interval $[i,j]$), and $P$ for the transposition (swap) operator
on $V \otimes V$.

%% =================================================================
\section{Part 1: Jimbo's theorem (classical)}
\label{sec:part1}
%% =================================================================

\begin{theorem}[Jimbo \cite{jimbo1986}]
\label{thm:jimbo}
For $n \ge k$, the map $T_i \mapsto \check{R}_i$ defines an algebra isomorphism
\[
  H_q(S_k) \;\xrightarrow{\;\sim\;}\; \End_{U_q(\mathfrak{sl}_n)}(V^{\otimes k}),
\]
where $\check{R}_i = \check{R}_{i,i+1}$ is the $R$-matrix braiding acting on
the $i$-th and $(i+1)$-th tensor factors.
\end{theorem}

The $R$-matrix on $V \otimes V$ takes the explicit form (in the basis
$\{e_a \otimes e_b\}$):
\begin{equation}\label{eq:R-matrix}
  \check{R}(e_a \otimes e_b) = \begin{cases}
    q\, e_a \otimes e_a & \text{if } a = b,\\
    e_b \otimes e_a + (q - q^{-1})\, e_a \otimes e_b & \text{if } a < b,\\
    e_b \otimes e_a & \text{if } a > b.
  \end{cases}
\end{equation}
This satisfies \eqref{eq:hecke-quad} and the Yang--Baxter equation on $V^{\otimes 3}$.

%% =================================================================
\section{Part 2: Cactus group at \texorpdfstring{$q=0$}{q=0}}
\label{sec:part2}
%% =================================================================

\begin{theorem}[Cactus representation at $q=0$]
\label{thm:cactus-q0}
At $q = 0$, the $R$-matrix degenerates to
\begin{equation}\label{eq:R-at-0}
  \check{R}\big|_{q=0}(e_a \otimes e_b) = \begin{cases}
    0 & \text{if } a = b,\\
    e_b \otimes e_a - e_a \otimes e_b & \text{if } a < b,\\
    e_b \otimes e_a & \text{if } a > b.
  \end{cases}
\end{equation}
The swap operator $P = 2\check{R}\big|_{q=0} + I$ satisfies the cactus relations:
defining $\sigma_{[i,j]} = w_0^{[i,j]}(P)$ (the longest element of $S_{\{i,\ldots,j\}}$
with each transposition $s_\ell$ realized by $P_{\ell,\ell+1}$),
the operators $\sigma_{[i,j]}$ satisfy \emph{(C1)--(C3)} on $V^{\otimes k}$.
\end{theorem}

\begin{proof}
\textbf{Step 1: $P$ from $\check{R}$.}
At $q=0$, the Hecke relation \eqref{eq:hecke-quad} becomes $T(T+1) = 0$,
so $T^2 = -T$. Setting $P = 2T + I$ (i.e.\ $\pi = T + 1$ in the Demazure notation),
we compute $P^2 = 4T^2 + 4T + I = -4T + 4T + I = I$.
Direct evaluation confirms that $P$ acts as the coordinate swap on $V \otimes V$:
$P(e_a \otimes e_b) = e_b \otimes e_a$.

\textbf{Step 2: $P$ commutes with $\GL_n$.}
Since $P$ is the standard flip, it commutes with $\Delta^{(0)}(g) = g \otimes g$
for all $g \in \GL_n(\CC)$. (At $q=0$, the quantum group degenerates and the
coproduct becomes cocommutative on the $\GL_n$ part.)

\textbf{Step 3: Cactus relations.}
Since $P_i = P_{i,i+1}$ is the standard swap, $\sigma_{[i,j]} = w_0^{[i,j]}(P)$
is the permutation matrix for the reversal of positions $\{i,\ldots,j\}$.
The cactus relations (C1)--(C3) reduce to identities about such reversals
in the symmetric group $S_k$, which hold by the standard presentation of $J_k$
mapping to $S_k$ (see Henriques--Kamnitzer \cite{HK2006}).

\textbf{Step 4: Factoring through $H_0$.}
At $q=0$, $\End_{U_0}(V^{\otimes k})$ is a quotient of $H_0(S_k)$ by Jimbo's theorem
(which holds at $q=0$ as a surjection for all $n \ge k$; injectivity follows from
dimension counting as in \cite{jimbo1986}).
Each $P_i = 2T_i + 1 = 2\pi_i - 1$ is an element of $H_0(S_k)$,
hence $\sigma_{[i,j]} = w_0^{[i,j]}(P) \in H_0(S_k)$.
\end{proof}

\begin{remark}
The Demazure operators $\pi_i = T_i + 1$ satisfy $\pi_i^2 = \pi_i$
(idempotent) and the braid relation $\pi_i \pi_{i+1} \pi_i = \pi_{i+1} \pi_i \pi_{i+1}$.
The cactus generators $\sigma_{[i,j]}$ are expressed as reduced words in the
$\pi_i$ (or equivalently in $P_i = 2\pi_i - 1$).
\end{remark}

\paragraph{Computational verification.}
Verified for $\mathfrak{sl}_n$ ($n = 2, 3$) with $k = 3, 4, 5$:
all cactus relations (C1)--(C3) hold exactly.
At generic $q \ne 0$, (C1) and (C3) hold but the nesting relation (C2) \emph{fails}
--- the failure is proportional to $(q - q^{-1})^2$, confirming that the cactus structure
is specific to the $q = 0$ degeneration when using swap operators.

%% =================================================================
\section{Part 3: Multiplicity bundle theorem (reference)}
\label{sec:part3}
%% =================================================================

The following was proved in the companion paper \cite{clio-mult-bundle}.

\begin{theorem}[Multiplicity Bundle Theorem]
\label{thm:mult-bundle}
Under Schur--Weyl duality $V^{\otimes k} \cong \bigoplus_\lambda V_\lambda \otimes S_\lambda$
(as $\GL_n \times S_k$-bimodule), the coboundary commutor decomposes as
\[
  \sigma_{[i,j]} = \bigoplus_\lambda \bigl(\id_{V_\lambda} \otimes \rho_{S_\lambda}(w_0^{[i,j]})\bigr).
\]
The crystal functor $F: \Rep(U_0) \to \Crys$ kills the multiplicity factor:
$F(\sigma_{[i,j]}) = \id$.
\end{theorem}

This establishes that the commutor lives entirely in the ``multiplicity bundle''
--- the $S_\lambda$ factor --- and is invisible to crystal combinatorics.

%% =================================================================
\section{Part 4: Coboundary structure in \texorpdfstring{$H_q$}{Hq} for all \texorpdfstring{$q$}{q}}
\label{sec:part4}
%% =================================================================

This is the key new result. We prove it by combining two known theorems.

\begin{theorem}[Coboundary commutor in $H_q$]
\label{thm:commutor-in-Hq}
For all $q$ (generic), the Kamnitzer--Tingley coboundary commutor
$\sigma_{V,W}: V \otimes W \to W \otimes V$ on finite-dimensional $U_q(\mathfrak{sl}_n)$-modules
restricts, on tensor powers of the fundamental representation, to elements of $H_q(S_k)$.
Specifically, for each interval $[i,j] \subseteq [1,k]$, the commutor
$\sigma_{[i,j]} \in \End(V^{\otimes k})$ lies in $H_q(S_k)$.
\end{theorem}

\begin{proof}
We combine two results:

\textbf{Theorem A} (Kamnitzer--Tingley \cite{KT2009}):
$\Rep(U_q(\mathfrak{g}))$ is a coboundary monoidal category.
The coboundary commutor $\sigma_{V,W}$ is defined via Drinfeld's unitarized $R$-matrix:
if $R_{V,W}: V \otimes W \to V \otimes W$ is the universal $R$-matrix and
$u$ is the Drinfeld element, then
\[
  \sigma_{V,W} = P_{V,W} \circ R_{V,W}^{\mathrm{uni}},
\]
where $R^{\mathrm{uni}}$ is the unitarized (self-adjoint) form of $R$.
By construction, $\sigma_{V,W}$ is a $U_q(\mathfrak{g})$-module isomorphism
$V \otimes W \xrightarrow{\sim} W \otimes V$, and the family $\{\sigma_{V,W}\}$
satisfies the coboundary (cactus) axiom.

In particular, for any subinterval $[i,j]$ and the associated tensor product
decomposition, the commutor $\sigma_{[i,j]}$ is a $U_q(\mathfrak{g})$-equivariant
endomorphism of $V^{\otimes k}$:
\[
  \sigma_{[i,j]} \in \End_{U_q(\mathfrak{sl}_n)}(V^{\otimes k}).
\]

\textbf{Theorem B} (Jimbo, Theorem~\ref{thm:jimbo}):
For $n \ge k$,
\[
  \End_{U_q(\mathfrak{sl}_n)}(V^{\otimes k}) \cong H_q(S_k).
\]

Combining: $\sigma_{[i,j]} \in \End_{U_q(\mathfrak{sl}_n)}(V^{\otimes k}) \cong H_q(S_k)$.

Since $\sigma_{[i,j]}$ satisfies the cactus relations (C1)--(C3)
(by the Kamnitzer--Tingley theorem --- this is the coboundary axiom for
the monoidal category), these relations hold \emph{inside $H_q(S_k)$}.
\end{proof}

\begin{remark}[The Drinfeld twist is essential]
\label{rmk:drinfeld}
The unitarized $R$-matrix $R^{\mathrm{uni}}$ involves a twist by the Drinfeld element $u$.
On $V \otimes V$, the normalized commutor takes the form
\[
  \sigma = \frac{2\check{R} - (q - q^{-1})I}{q + q^{-1}},
\]
which satisfies $\sigma^2 = I$ (involutivity) and is a $U_q$-intertwiner,
but does \emph{not} satisfy the braid relation. The failure is:
\[
  \sigma_1 \sigma_2 \sigma_1 - \sigma_2 \sigma_1 \sigma_2 = -\frac{2(q-q^{-1})^2}{(q+q^{-1})^3}(T_1 - T_2).
\]
This is not a contradiction: the cactus relations for $\sigma_{[i,j]}$ on
\emph{higher tensor products} involve the Drinfeld twist applied to composite
subsystems, which cannot be factored as products of pairwise commutors.
The key point is that the KT construction produces the correct commutor
\emph{directly} at the level of each subinterval, bypassing the need to
compose pairwise braidings.
\end{remark}

\begin{remark}[What $H_q$ simultaneously contains]
At generic $q$, the Hecke algebra $H_q(S_k)$ contains:
\begin{itemize}
  \item \textbf{Braiding data}: the generators $T_i$ (images of the $R$-matrix),
    satisfying the braid relation \eqref{eq:braid} and generating a braid group
    representation $B_k \to H_q(S_k)^\times$.
  \item \textbf{Coboundary data}: the Kamnitzer--Tingley commutors $\sigma_{[i,j]}$,
    satisfying the cactus relations and generating a cactus group representation
    $J_k \to H_q(S_k)^\times$.
\end{itemize}
Both structures coexist for all $q$. At $q=0$, the braiding degenerates
($\check{R}$ becomes nilpotent: $\check{R}^2 = -\check{R}$),
but the coboundary structure simplifies to the swap-based cactus action
of Part~2.
\end{remark}

\paragraph{Computational verification.}
On $V^{\otimes 2}$ ($\mathfrak{sl}_2$): $\sigma^2 = I$ verified for symbolic $q$.
The commutor $\sigma = (2\check{R}-(q-q^{-1})I)/(q+q^{-1})$ is a $U_q$-intertwiner
(commutes with $\Delta(E), \Delta(F), \Delta(K)$) for all $q$.
The swap $P$ is \emph{not} a $U_q$-intertwiner for $q \ne 1$:
$[P, \Delta(E)] \ne 0$.
This confirms that the Drinfeld twist genuinely alters the commutor.

%% =================================================================
\section{Part 5: Semisimplification and crystals}
\label{sec:part5}
%% =================================================================

\begin{theorem}[Algebra chain mirrors category chain]
\label{thm:part5}
The three-level categorical hierarchy
\[
  \Rep(U_q(\mathfrak{sl}_n)) \;\xrightarrow{q \to 0}\; \Rep(U_0(\mathfrak{sl}_n))
  \;\xrightarrow{F}\; \Crys(\mathfrak{sl}_n)
\]
is mirrored by the algebra chain
\[
  H_q(S_k) \;\xrightarrow{q \to 0}\; H_0(S_k) \;\xrightarrow{\mathrm{ss}}\; \CC[S_k],
\]
where $\mathrm{ss}$ denotes the semisimplification map $T_w \mapsto w$.
At the character level, this passage is recorded by the Demazure atom decomposition:
\begin{equation}\label{eq:atom-decomp}
  s_\lambda = \sum_{\substack{\alpha \in \ZZ_{\ge 0}^n \\ \alpha_+ = \lambda}} A_\alpha,
\end{equation}
where $A_\alpha$ is the Demazure atom (alternating Demazure operator) and $\alpha_+$
is the dominant rearrangement of $\alpha$.
\end{theorem}

\begin{proof}
The proof proceeds in three steps, establishing the algebra chain, its representation-theoretic
content, and the character-level identification.

\textbf{Step 1: The algebra chain.}
The specialization $H_q(S_k) \xrightarrow{q \to 0} H_0(S_k)$ is literal:
set $q = 0$ in the Hecke relation \eqref{eq:hecke-quad} to obtain $T_i(T_i + 1) = 0$.
The map $H_0(S_k) \xrightarrow{\mathrm{ss}} \CC[S_k]$ sends $T_w \mapsto w$;
this is well-defined because the braid and commutation relations
\eqref{eq:braid}--\eqref{eq:commute} hold in both algebras, and the quadratic
relation $T^2 = -T$ in $H_0$ maps to $s^2 = 1$ in $\CC[S_k]$
(the image of $T^2 + T = 0$ under $T \mapsto s$ is $s^2 + s = s(s+1) = 0$,
which does \emph{not} hold in $\CC[S_k]$ --- hence this is not an algebra
homomorphism but rather the semisimplification at the module level, as we now explain).

More precisely: $H_0(S_k)$ is \emph{not} semisimple (Norton \cite{norton1979}).
It has $2^{k-1}$ simple modules, one for each composition of $k$
(equivalently, one for each subset $D \subseteq \{1,\ldots,k-1\}$ of descents).
Each simple module is $1$-dimensional, with $\pi_i = T_i + 1$ acting as $1$ if
$i \in D$ and as $0$ if $i \notin D$.

The group algebra $\CC[S_k]$ is semisimple with $p(k)$ simple modules
(Specht modules $S_\lambda$, one per partition $\lambda \vdash k$).

The semisimplification is the functor $H_0\Mod \to \CC[S_k]\Mod$ induced by the
$\CC$-linear map sending $T_w \mapsto w$ and extending by restriction of scalars.
This merges composition-indexed modules into partition-indexed ones: the $H_0$-modules
indexed by compositions that rearrange to the same partition $\lambda$ are identified
with the single Specht module $S_\lambda$.

\textbf{Step 2: Matching the categorical chain.}
Under Jimbo's isomorphism (Theorem~\ref{thm:jimbo}):
\begin{itemize}
  \item $\Rep(U_q(\mathfrak{sl}_n))$ on $V^{\otimes k}$: the endomorphism algebra is $H_q(S_k)$.
  \item $\Rep(U_0(\mathfrak{sl}_n))$ on $V^{\otimes k}$: the endomorphism algebra is $H_0(S_k)$.
    The non-semisimplicity of $H_0$ reflects the fact that crystal bases can
    distinguish copies of the same irreducible that are related by different
    highest-weight paths.
  \item $\Crys(\mathfrak{sl}_n)$ on $B(\omega_1)^{\otimes k}$: the crystal limit
    forgets the $H_0$-action (by the Multiplicity Bundle Theorem~\ref{thm:mult-bundle},
    the commutor acts trivially on crystals). The residual symmetry is $S_k$,
    which permutes tensor factors via $P$ (the swap).
\end{itemize}

For example, at $k=3$, $n \ge 3$: $H_0(S_3)$ has 4 simple modules
(indexed by compositions $3, 21, 12, 111$). The Schur--Weyl decomposition
$V^{\otimes 3} \cong V_{(3)} \otimes S_{(3)} \oplus V_{(2,1)} \otimes S_{(2,1)}
\oplus V_{(1,1,1)} \otimes S_{(1,1,1)}$ has two copies of $V_{(2,1)}$
(one for each standard Young tableau of shape $(2,1)$). These two copies
are distinguished by $H_0$ (via descent sets $\{1\}$ vs $\{2\}$) but identified
by $\CC[S_3]$ (both live in the single Specht module $S_{(2,1)}$).

\textbf{Step 3: Character-level identification via Demazure atoms.}
The Demazure operator $\pi_i$ acts on polynomials $f \in \CC[x_1,\ldots,x_n]$ by
\[
  \pi_i(f) = \frac{x_i f - x_{i+1}\, s_i(f)}{x_i - x_{i+1}},
\]
where $s_i$ swaps $x_i$ and $x_{i+1}$.
The \emph{key polynomial} (Demazure character) is
$\kappa_\alpha = \pi_{i_1} \cdots \pi_{i_r}(x^{\alpha_+})$
for a reduced word $s_{i_1} \cdots s_{i_r}$ taking $\alpha_+$ to $\alpha$.
The \emph{Demazure atom} is
$A_\alpha = \pi_{i_1}^{(0)} \cdots \pi_{i_r}^{(0)}(x^{\alpha_+})$,
where $\pi_i^{(0)} = \pi_i - \id$.

The identity \eqref{eq:atom-decomp} is the character-level shadow of the
semisimplification functor: each atom $A_\alpha$ corresponds to a $1$-dimensional
$H_0$-module (indexed by the descent set of $\alpha$), and summing over all
compositions with $\alpha_+ = \lambda$ recovers the Schur function $s_\lambda$
--- the character of the $\CC[S_k]$-module $S_\lambda$ (or rather, the $\GL_n$-character
of $V_\lambda$). The map $A_\alpha \mapsto s_{\alpha_+}$ is precisely
the character of the semisimplification $H_0\Mod \to \CC[S_k]\Mod$.

This was observed by Lascoux--Sch\"utzenberger \cite{LS1990} (for key polynomials)
and made precise in the $H_0$ framework by Mason \cite{mason2009} and
Pun--Woodruff (see also Haglund--Luoto--Mason--van~Willigenburg \cite{HLMW2011}).
\end{proof}

\paragraph{Computational verification.}
Identity \eqref{eq:atom-decomp} verified for $n = 3$ and all partitions $\lambda$
with $|\lambda| \le 5$: specifically
$\lambda = (3), (2{,}1), (1{,}1{,}1), (2{,}2), (2{,}1{,}1), (3{,}1), (4), (3{,}2)$.
In each case, $s_\lambda = \sum_{\alpha_+ = \lambda} A_\alpha$ holds exactly.
Key polynomials $\kappa_\alpha$ are verified to have non-negative integer coefficients
and to be nested: $\kappa_\alpha = \sum_{\beta \le \alpha} A_\beta$ where $\le$ is the
weak Bruhat order.

%% =================================================================
\section{The Hecke Transition Algebra Theorem (combined)}
\label{sec:combined}
%% =================================================================

Combining Parts 1--5:

\begin{theorem}[Hecke Transition Algebra]
\label{thm:main}
Let $V = \CC^n$ be the fundamental representation of $U_q(\mathfrak{sl}_n)$, $n \ge k$.
\begin{enumerate}
  \item \textbf{(Jimbo)} $\End_{U_q(\mathfrak{sl}_n)}(V^{\otimes k}) \cong H_q(S_k)$.
  \item \textbf{(Cactus at $q=0$)} The swap operators $P_i = 2T_i + 1\big|_{q=0}$
    define a cactus group representation $J_k \to H_0(S_k)^\times$.
  \item \textbf{(Multiplicity bundle)} The commutor decomposes as
    $\sigma_{[i,j]} = \bigoplus_\lambda (\id_{V_\lambda} \otimes \rho_{S_\lambda}(w_0^{[i,j]}))$
    and is killed by the crystal functor.
  \item \textbf{(Coboundary in $H_q$)} For all generic $q$, the Kamnitzer--Tingley
    commutors $\sigma_{[i,j]}$ lie in $H_q(S_k)$ and satisfy the cactus relations therein.
  \item \textbf{(Semisimplification = crystal)} The chain
    $H_q \to H_0 \to \CC[S_k]$ mirrors
    $\Rep(U_q) \to \Rep(U_0) \to \Crys$,
    with the passage $H_0 \to \CC[S_k]$ corresponding to
    forgetting the Demazure atom decomposition $s_\lambda = \sum A_\alpha$.
\end{enumerate}
The Hecke algebra $H_q(S_k)$ is the \emph{transition algebra}: it simultaneously
carries the braiding (braid group via $T_i$) and coboundary (cactus group via
$\sigma_{[i,j]}$) structures, with the parameter $q$ governing the transition
between them.
\end{theorem}

%% =================================================================
\section{Gaps and open directions}
\label{sec:gaps}
%% =================================================================

\begin{enumerate}
  \item \textbf{Part 4, explicit elements.}
  The proof of Theorem~\ref{thm:commutor-in-Hq} is by abstract argument
  (KT commutor is a $U_q$-intertwiner, hence in $H_q$ by Jimbo).
  Writing $\sigma_{[i,j]}$ as an \emph{explicit} element of $H_q(S_k)$
  (i.e., as a polynomial in the $T_i$) remains open for $k \ge 3$.
  On $V^{\otimes 2}$, $\sigma = (2T - (q-q^{-1}))/(q+q^{-1})$.
  On $V^{\otimes 3}$, $\sigma_{[1,3]}$ involves the Drinfeld twist on
  the composite system and does not factor as a product of pairwise terms.

  \item \textbf{Part 5, functoriality.}
  The semisimplification $H_0\Mod \to \CC[S_k]\Mod$ should be identified
  precisely with the crystal functor $F$. We have verified this at the level
  of characters (Demazure atoms $\to$ Schur) and on specific modules ($k \le 5$),
  but a complete proof of the equivalence of functors would require showing that
  $F$ acts on morphisms as $H_0$-module restriction. This should follow from
  the Multiplicity Bundle Theorem but deserves a careful write-up.

  \item \textbf{Connection to Kalmykov.}
  The double affine Hecke algebra (DAHA) contains $H_q$ and is isomorphic
  to the shuffle algebra. The transition algebra therefore lives inside
  the shuffle substrate, suggesting connections to the KPZ universality
  class and integrable probability.
\end{enumerate}

%% =================================================================
\section*{References}
%% =================================================================

\begin{thebibliography}{99}
\bibitem{jimbo1986}
M.~Jimbo, \emph{A $q$-analogue of $U(\mathfrak{gl}(N+1))$, Hecke algebra,
and the Yang--Baxter equation}, Lett.\ Math.\ Phys.\ \textbf{11} (1986), 247--252.

\bibitem{KT2009}
J.~Kamnitzer and P.~Tingley, \emph{The crystal commutor and Drinfeld's unitarized
$R$-matrix}, J.\ Algebraic Combin.\ \textbf{29} (2009), 315--335.
arXiv:0707.2248.

\bibitem{HK2006}
A.~Henriques and J.~Kamnitzer, \emph{Crystals and coboundary categories},
Duke Math.\ J.\ \textbf{132} (2006), 191--216.

\bibitem{norton1979}
P.~N.~Norton, \emph{$0$-Hecke algebras}, J.\ Austral.\ Math.\ Soc.\ Ser.~A
\textbf{27} (1979), 337--357.

\bibitem{LS1990}
A.~Lascoux and M.-P.~Sch\"utzenberger, \emph{Keys and standard bases},
Invariant Theory and Tableaux (Minneapolis, 1988), IMA Vol.\ Math.\ Appl.\ \textbf{19},
Springer, 1990, pp.\ 125--144.

\bibitem{mason2009}
S.~Mason, \emph{An explicit construction of type A Demazure atoms},
J.\ Algebraic Combin.\ \textbf{29} (2009), 295--313.

\bibitem{HLMW2011}
J.~Haglund, K.~Luoto, S.~Mason, and S.~van~Willigenburg,
\emph{Quasisymmetric Schur functions}, J.\ Combin.\ Theory Ser.~A
\textbf{118} (2011), 463--490.

\bibitem{clio-mult-bundle}
Clio, \emph{The Multiplicity Bundle Theorem}, working paper, April 11, 2026.

\end{thebibliography}

\end{document}
